% !TITLE 赤壁
% !AUTHOR 杜牧
% !DYNASTY 唐
% !GENRE 七言绝句
% !ORDER 57

\begin{poem}
折戟沉沙铁未销\textsuperscript{1}，自将磨洗认前朝\textsuperscript{2}。
东风不与周郎便\textsuperscript{3}，铜雀春深锁二乔\textsuperscript{4}。
\end{poem}

\begin{annotations}
\notetext{1}{折戟沉沙铁未销：折断的战戟沉埋在沙中，铁还没有完全锈蚀。戟（jǐ），古代一种兵器。}
\notetext{2}{自将磨洗认前朝：自己拿起来磨去锈迹、洗去泥沙，认出是前朝（三国）赤壁之战的遗物。}
\notetext{3}{周郎：指周瑜（175-210），当时人称“周郎”。赤壁之战中用火攻大破曹操，据说天时起了关建作用——东风助火势烧向曹军。}
\notetext{4}{铜雀春深锁二乔：铜雀台是曹操晚年享乐的宫殿。二乔是大乔（嫁孙策）和小乔（嫁周瑜）。杜牧假设：如果没有东风，周瑜战败，二乔就会被曹操关在铜雀台中——用两个美人的命运来议论历史的偶然性。}
\end{annotations}

\begin{appreciation}
这是杜牧的一首咏史诗，也是中国文学史上谈历史的偶然性最出名的一首之一。杜牧在赤壁的江边捡到一支断戟，六百多年前赤壁之战的遗物，由此发了一通议论。

折戟沉沙铁未销——戟是一种兵器，一头刺，一头钩。沉沙，说明它在江底躺了很久；铁未销，泡了六百多年还没锈烂——既说明当年的铁打得好，也像那段历史一直没被忘记。

自将磨洗认前朝——将，拿。磨、洗，两个动作：一个晚唐的诗人，蹲在江边，把一支三国时代的断戟擦洗干净，像和六百年前的历史握了个手。

东风不与周郎便——便，成全、方便。赤壁之战，周瑜用火攻，靠的是一阵东南风。杜牧说：要是没有那阵风呢？

铜雀春深锁二乔——没那阵风，江东就完了，大乔小乔要被关进曹操的铜雀台。铜雀台是曹操在邺城盖的，晚年享乐的地方。二乔，大乔嫁孙策，小乔嫁周瑜。锁字很冷，春深两个字又暖又深——最好的时节，被锁在最高的台子里。

有人批评过这种写法轻佻，但杜牧的高明恰在这里：历史的重量，落在具体的人身上才算数。不谈帝业兴衰，谈两个人会不会被关进一座台子——这是把历史从年表里捞出来，放到人身上。

这本书后面有苏轼的《念奴娇·赤壁怀古》，那里写小乔初嫁了，雄姿英发——苏轼写的是小乔的幸福，杜牧假设的是小乔的不幸，同一个女子，一边是现实，一边是假设。还有曹操自己，书里《观沧海》，东临碣石，以观沧海——那个站在碣石上看海、雄心万丈的曹操，一年之后就在赤壁吃了这股东风的亏。两首诗拼在一起，曹操的一生就齐了。

历史常常系于一阵风。但风来之前，船是提前造好的，周瑜的火攻船，是他自己备下的。你将来考试也好，做事也好，别指望东风，把该磨的戟磨好。哪天捡到什么旧东西，也拿起来磨一磨、洗一洗，看看是哪个朝代的。
\end{appreciation}
